\documentclass[12pt,oneside]{book}

\usepackage[spanish,es-nodecimaldot]{babel}
\usepackage[utf8]{inputenc}
\usepackage[T1]{fontenc}
\usepackage{lmodern}
\usepackage[a4paper,margin=2.5cm]{geometry}
\usepackage{microtype}
\usepackage{setspace}
\usepackage{amsmath,amssymb,amsthm}
\usepackage{physics}
\usepackage{siunitx}
\usepackage{graphicx}
\usepackage{xcolor}
\usepackage{enumitem}
\usepackage{hyperref}
\usepackage{bookmark}
\usepackage{natbib}

\hypersetup{
  colorlinks=true,
  linkcolor=blue,
  urlcolor=blue,
  citecolor=blue,
  pdftitle={Clase extensa: cosmología y modelado astrofísico contemporáneo},
  pdfauthor={Semillero de Astronomía}
}

\setstretch{1.15}
\setlength{\parskip}{0.45em}
\setlength{\parindent}{0pt}
\setcounter{secnumdepth}{3}
\setcounter{tocdepth}{2}

\newtheorem{definicion}{Definición}[chapter]
\newtheorem{ejemplo}{Ejemplo}[chapter]
\newtheorem{actividad}{Actividad}[chapter]
\newtheorem{observacion}{Observación}[chapter]

\newcommand{\LCDM}{$\Lambda$CDM}
\newcommand{\Hcero}{H_0}
\newcommand{\OmM}{\Omega_{\mathrm{m}}}
\newcommand{\OmL}{\Omega_{\Lambda}}
\newcommand{\OmR}{\Omega_{\mathrm{r}}}
\newcommand{\Omk}{\Omega_{k}}

\title{Clase Extensa\\[0.5em]
Cosmología Física, Inferencia de Datos y Lectura Crítica de Investigación Actual}
\author{Semillero de Astronomía -- IEEOH}
\date{Mayo de 2026}

\begin{document}

\frontmatter
\maketitle
\tableofcontents

\mainmatter

\chapter{Propósito de la clase y artículos base}

Esta clase fue diseñada para estudiar, con profundidad, una ruta de trabajo que conecta
la física teórica, el análisis cuantitativo y la interpretación crítica de literatura científica.
La base principal de lectura son los tres artículos solicitados:
\begin{itemize}[leftmargin=1.5em]
  \item arXiv:2603.02372v1 \citep{arxiv2603};
  \item arXiv:2001.11644 \citep{arxiv2001};
  \item Scientific Reports (Nature Portfolio), artículo \texttt{s41598-024-54700-x} \citep{srep2024}.
\end{itemize}

Además, para consolidar el marco teórico y matemático se incorporan textos de apoyo estándar
en cosmología, estadística e instrumentación astronómica \citep{dodelson,ryden,hogg1999,trotta2008,wall2012}.

\section{Objetivos formativos}
Al finalizar esta clase, el estudiante deberá ser capaz de:
\begin{enumerate}[leftmargin=1.5em]
  \item explicar el marco cosmológico homogéneo e isótropo y su ecuación dinámica básica;
  \item usar relaciones matemáticas para conectar teoría con observables astronómicos;
  \item resolver problemas cuantitativos con unidades, propagación de incertidumbre y ajuste de modelos;
  \item identificar supuestos, limitaciones y alcances de artículos científicos contemporáneos;
  \item proponer actividades experimentales y de simulación para profundizar en el tema.
\end{enumerate}

\section{Ruta pedagógica}
La clase se divide en cuatro bloques:
\begin{enumerate}[leftmargin=1.5em]
  \item Marco teórico físico (geometría, expansión y contenido del universo).
  \item Formalismo matemático (distancias, crecimiento y modelos efectivos).
  \item Metodología de datos (estimación, inferencia y validación).
  \item Taller de ejemplos extensos y tareas finales.
\end{enumerate}

\chapter{Marco teórico: de la geometría cósmica al universo observado}

\section{Principio cosmológico y métrica FLRW}

El \textbf{principio cosmológico} afirma que, a gran escala, el universo es homogéneo e isótropo.
Bajo este supuesto, la geometría espacio-temporal se modela con la métrica de
Friedmann--Lemaître--Robertson--Walker (FLRW):
\begin{equation}
  ds^2 = -c^2dt^2 + a^2(t)\left[\frac{dr^2}{1-kr^2}+r^2(d\theta^2+\sin^2\theta\,d\phi^2)\right],
\end{equation}
donde $a(t)$ es el factor de escala, $k\in\{-1,0,+1\}$ codifica la curvatura espacial y
$c$ es la velocidad de la luz.

\begin{definicion}[Parámetro de Hubble]
Se define como
\begin{equation}
  H(t)=\frac{\dot a(t)}{a(t)},
\end{equation}
y su valor actual es $H(t_0)=\Hcero$.
\end{definicion}

\section{Ecuaciones de Friedmann y contenido energético}

A partir de la relatividad general, para un fluido perfecto se obtienen:
\begin{align}
  H^2 &= \frac{8\pi G}{3}\rho -\frac{kc^2}{a^2},\\
  \frac{\ddot a}{a} &= -\frac{4\pi G}{3}\left(\rho+\frac{3p}{c^2}\right).
\end{align}

La ecuación de conservación energía-momento en el fondo homogéneo toma la forma:
\begin{equation}
  \dot\rho + 3H\left(\rho+\frac{p}{c^2}\right)=0.
\end{equation}

Si $p=w\rho c^2$, entonces:
\begin{equation}
  \rho(a)\propto a^{-3(1+w)}.
\end{equation}

Casos relevantes:
\begin{itemize}[leftmargin=1.5em]
  \item materia no relativista ($w=0$): $\rho_m\propto a^{-3}$;
  \item radiación ($w=1/3$): $\rho_r\propto a^{-4}$;
  \item energía oscura tipo constante cosmológica ($w=-1$): $\rho_\Lambda=\text{cte}$.
\end{itemize}

\section{Forma adimensional y expansión con corrimiento al rojo}

Usando $1+z=1/a$ (con $a_0=1$), para \LCDM\ se tiene:
\begin{equation}
  E^2(z)\equiv\frac{H^2(z)}{H_0^2}=
  \OmR(1+z)^4+\OmM(1+z)^3+\Omk(1+z)^2+\OmL.
\end{equation}

En un universo plano ($\Omk=0$):
\begin{equation}
  \OmR+\OmM+\OmL=1.
\end{equation}

\begin{observacion}
La interpretación física de $E(z)$ es central: cada término domina en una época distinta
(radiación, materia, aceleración tardía). Esta idea aparece recurrentemente en literatura
cosmológica moderna y es indispensable para comparar teoría y observación
\citep{dodelson,ryden,arxiv2603,arxiv2001,srep2024}.
\end{observacion}

\section{Distancias cosmológicas observables}

Para una fuente a redshift $z$, la distancia comóvil radial en universo plano es:
\begin{equation}
  D_C(z)=c\int_0^z\frac{dz'}{H(z')}.
\end{equation}

Distancia de luminosidad y distancia angular:
\begin{align}
  D_L(z)&=(1+z)D_C(z),\\
  D_A(z)&=\frac{D_C(z)}{1+z}=\frac{D_L(z)}{(1+z)^2}.
\end{align}

Estas relaciones conectan observaciones fotométricas y geometría del universo,
fundamentales en análisis astrofísicos contemporáneos \citep{hogg1999}.

\chapter{Marco matemático de modelado e inferencia}

\section{Aproximación de bajo redshift}

Para $z\ll1$, una expansión de Taylor permite aproximar:
\begin{equation}
  D_L(z)\approx\frac{cz}{H_0}\left[1+\frac{1-q_0}{2}z\right],
\end{equation}
donde $q_0=-\ddot a a/\dot a^2\vert_{t_0}$ es el parámetro de desaceleración actual.

Esta aproximación es útil para ejercicios didácticos y chequeos de consistencia.

\section{Problema directo e inverso}

\textbf{Problema directo:} dados parámetros cosmológicos $\Theta=(H_0,\OmM,\OmL,\dots)$,
calcular observables esperados $\mathcal{O}_{\text{th}}(z;\Theta)$.

\textbf{Problema inverso:} a partir de datos $\mathcal{D}$, estimar los parámetros $\Theta$.

La función de verosimilitud gaussiana típica:
\begin{equation}
  \mathcal{L}(\Theta)\propto \exp\left[-\frac{1}{2}\chi^2(\Theta)\right],
\end{equation}
con
\begin{equation}
  \chi^2(\Theta)=\sum_{i=1}^{N}\frac{\left[\mathcal{O}_i-\mathcal{O}_{\text{th}}(z_i;\Theta)\right]^2}{\sigma_i^2}.
\end{equation}

\section{Inferencia bayesiana}

La inferencia moderna usa:
\begin{equation}
  P(\Theta\mid\mathcal{D})=\frac{\mathcal{L}(\Theta)\,\pi(\Theta)}{\mathcal{Z}},
\end{equation}
donde $\pi(\Theta)$ es la priori y $\mathcal{Z}$ la evidencia bayesiana.

Para comparación de modelos $M_1,M_2$:
\begin{equation}
  B_{12}=\frac{\mathcal{Z}_1}{\mathcal{Z}_2}.
\end{equation}

\section{Incertidumbre, sesgo y robustez}

Cualquier conclusión física debe separar:
\begin{itemize}[leftmargin=1.5em]
  \item \textbf{error estadístico:} dispersión por tamaño muestral y ruido;
  \item \textbf{error sistemático:} calibración, selección, modelo incompleto;
  \item \textbf{sesgo de interpretación:} sobreadaptación o uso de supuestos no verificados.
\end{itemize}

La lectura crítica de artículos, incluyendo los tres textos base de esta clase,
requiere identificar explícitamente estos tres niveles \citep{trotta2008,wall2012}.

\chapter{Ejemplos extensos resueltos}

\section{Ejemplo 1: evolución de densidades relativas con el redshift}

\begin{ejemplo}[Escalamiento de materia y radiación]
Suponga que hoy se tiene $\OmM=0.30$ y $\OmR=9\times10^{-5}$. Halle el redshift
aproximado de igualdad materia-radiación.

\textbf{Solución.}
La igualdad ocurre cuando
\begin{equation}
  \rho_m(z_{\mathrm{eq}})=\rho_r(z_{\mathrm{eq}}).
\end{equation}
Usando los escalados:
\begin{equation}
  \rho_{m0}(1+z_{\mathrm{eq}})^3=\rho_{r0}(1+z_{\mathrm{eq}})^4
  \Rightarrow 1+z_{\mathrm{eq}}=\frac{\rho_{m0}}{\rho_{r0}}=\frac{\OmM}{\OmR}.
\end{equation}
Entonces:
\begin{equation}
  1+z_{\mathrm{eq}}\approx\frac{0.30}{9\times10^{-5}}\approx3333.3
  \Rightarrow z_{\mathrm{eq}}\approx3332.
\end{equation}
Resultado físicamente consistente con escalas cosmológicas estándar.
\end{ejemplo}

\section{Ejemplo 2: cálculo de distancia de luminosidad a bajo redshift}

\begin{ejemplo}[Uso de expansión de Taylor]
Considere $H_0=70\,\si{km.s^{-1}.Mpc^{-1}}$, $q_0=-0.55$ y $z=0.10$.
Aproxime $D_L$ en Mpc.

\textbf{Solución.}
\begin{equation}
  D_L\approx\frac{cz}{H_0}\left[1+\frac{1-q_0}{2}z\right].
\end{equation}
Primero,
\begin{equation}
  \frac{cz}{H_0}=\frac{(3\times10^5\,\si{km/s})(0.10)}{70\,\si{km.s^{-1}.Mpc^{-1}}}
  \approx428.57\,\si{Mpc}.
\end{equation}
Factor correctivo:
\begin{equation}
  1+\frac{1-(-0.55)}{2}(0.10)=1+\frac{1.55}{2}(0.10)=1+0.0775=1.0775.
\end{equation}
Finalmente,
\begin{equation}
  D_L\approx428.57\times1.0775\approx461.8\,\si{Mpc}.
\end{equation}
\end{ejemplo}

\section{Ejemplo 3: módulo de distancia y magnitud absoluta}

\begin{ejemplo}[Conversión observacional]
Si una fuente presenta magnitud aparente $m=18.6$ y se estima
$D_L=461.8\,\si{Mpc}$, encuentre la magnitud absoluta $M$.

\textbf{Solución.}
Usamos
\begin{equation}
  \mu = m-M = 5\log_{10}\left(\frac{D_L}{10\,\mathrm{pc}}\right).
\end{equation}
Como $D_L=461.8\times10^6\,\mathrm{pc}$:
\begin{equation}
  \mu=5\log_{10}(4.618\times10^7)=5(7+\log_{10}4.618)
  \approx5(7+0.664)=38.32.
\end{equation}
Por tanto:
\begin{equation}
  M=m-\mu=18.6-38.32=-19.72.
\end{equation}
Este valor es compatible con objetos transitorios muy luminosos.
\end{ejemplo}

\section{Ejemplo 4: ajuste por $\chi^2$ de un parámetro efectivo}

\begin{ejemplo}[Estimación de $H_0$ simplificada]
Suponga tres mediciones de distancia de luminosidad a bajo $z$:
\begin{center}
\begin{tabular}{c c c}
$z_i$ & $D_{L,i}$ (Mpc) & $\sigma_i$ (Mpc)\\
0.03 & 130 & 8\\
0.05 & 215 & 10\\
0.08 & 345 & 14
\end{tabular}
\end{center}
Con aproximación lineal $D_L\approx cz/H_0$, estime $H_0$ por mínimos cuadrados.

\textbf{Solución.}
Modelo lineal en la variable $x_i=cz_i$:
\begin{equation}
  D_{L,i}=\beta x_i,\qquad \beta=\frac{1}{H_0}.
\end{equation}
Estimador ponderado:
\begin{equation}
  \hat\beta=\frac{\sum_i w_i x_i D_{L,i}}{\sum_i w_i x_i^2},\qquad w_i=\frac{1}{\sigma_i^2}.
\end{equation}
Con $c=3\times10^5\,\si{km/s}$,
\begin{equation}
  x=(9000,15000,24000)\,\si{km/s}.
\end{equation}
Sustituyendo, se obtiene aproximadamente
\begin{equation}
  \hat\beta\approx1.43\times10^{-2}\,\si{Mpc/(km.s^{-1})}
  \Rightarrow \hat H_0\approx69.9\,\si{km.s^{-1}.Mpc^{-1}}.
\end{equation}
\end{ejemplo}

\section{Ejemplo 5: propagación de incertidumbre en función cosmológica}

\begin{ejemplo}[Error en $E(z)$]
Para universo plano materia+energía oscura:
\begin{equation}
  E(z)=\sqrt{\OmM(1+z)^3+(1-\OmM)}.
\end{equation}
En $z=1$, con $\OmM=0.30\pm0.02$, estime $\sigma_E$.

\textbf{Solución.}
\begin{equation}
  E(1)=\sqrt{0.30\times8+0.70}=\sqrt{3.1}=1.7607.
\end{equation}
Derivada parcial:
\begin{equation}
  \pdv{E}{\OmM}=\frac{(1+z)^3-1}{2E}.
\end{equation}
En $z=1$:
\begin{equation}
  \pdv{E}{\OmM}=\frac{8-1}{2(1.7607)}=\frac{7}{3.5214}\approx1.988.
\end{equation}
Error propagado:
\begin{equation}
  \sigma_E\approx\left|\pdv{E}{\OmM}\right|\sigma_{\OmM}
  \approx1.988\times0.02\approx0.040.
\end{equation}
Así, $E(1)=1.76\pm0.04$.
\end{ejemplo}

\section{Ejemplo 6: criterio de información para comparar modelos}

\begin{ejemplo}[AIC y penalización de complejidad]
Modelo A: $k_A=2$ parámetros, $\chi^2_{\min,A}=112$.\
Modelo B: $k_B=4$ parámetros, $\chi^2_{\min,B}=108$.
Compare con
\begin{equation}
  \mathrm{AIC}=\chi^2_{\min}+2k.
\end{equation}

\textbf{Solución.}
\begin{align}
  \mathrm{AIC}_A &= 112 + 2(2)=116,\\
  \mathrm{AIC}_B &= 108 + 2(4)=116.
\end{align}
Ambos modelos quedan empatados en este criterio; la mejora en ajuste de B
se compensa exactamente por su mayor complejidad.
\end{ejemplo}

\chapter{Lectura crítica guiada de los tres artículos base}

\section{Cómo leer investigación avanzada sin perder rigor}

En esta sección se propone un protocolo útil para seminarios:
\begin{enumerate}[leftmargin=1.5em]
  \item \textbf{Identificar la pregunta científica central.}
  \item \textbf{Aislar el modelo físico y sus supuestos.}
  \item \textbf{Escribir las ecuaciones clave en notación propia.}
  \item \textbf{Distinguir resultados robustos de resultados dependientes de calibración.}
  \item \textbf{Comparar con literatura previa y con datos independientes.}
\end{enumerate}

\section{Matriz de análisis para discusión en clase}

Para cada artículo base \citep{arxiv2603,arxiv2001,srep2024}, diligenciar:
\begin{center}
\begin{tabular}{p{0.24\textwidth} p{0.68\textwidth}}
\textbf{Pregunta central} & ¿Qué problema intenta resolver el artículo? \\
\textbf{Marco teórico} & ¿Qué ecuaciones y aproximaciones utiliza? \\
\textbf{Datos o simulaciones} & ¿Qué tipo de evidencia respalda las conclusiones? \\
\textbf{Limitaciones} & ¿Qué suposiciones podrían cambiar el resultado? \\
\textbf{Siguiente paso} & ¿Qué experimento/observación pondría a prueba la hipótesis? \\
\end{tabular}
\end{center}

\section{Conexión con el trabajo del semillero}

La meta pedagógica no es repetir pasivamente un artículo, sino aprender a:
\begin{itemize}[leftmargin=1.5em]
  \item traducir un resultado especializado a conceptos de bachillerato avanzado;
  \item validar números mediante cálculos independientes de orden de magnitud;
  \item reconocer cuándo un resultado depende críticamente del método estadístico;
  \item formular preguntas nuevas con base en evidencia cuantitativa.
\end{itemize}

\chapter{Actividades y tareas finales}

\section{Actividades en clase (trabajo colaborativo)}

\begin{actividad}[Actividad 1: Mapa conceptual físico-matemático]
En grupos de tres, construyan un mapa que conecte: métrica FLRW, ecuaciones de
Friedmann, función $E(z)$, distancias cosmológicas e inferencia estadística. Cada
flecha debe incluir una ecuación o una justificación matemática breve.
\end{actividad}

\begin{actividad}[Actividad 2: Debate de supuestos]
Seleccionen un supuesto fuerte de cada artículo base y discutan: (i) por qué se adopta,
(ii) qué ventaja aporta, y (iii) cómo podría sesgar la conclusión si no se cumple.
\end{actividad}

\begin{actividad}[Actividad 3: Reproducción parcial de resultados]
Elijan una figura o tabla numérica de uno de los artículos base e intenten reproducir,
con una hoja de cálculo o Python, una versión simplificada del resultado.
Documenten diferencias y posibles causas.
\end{actividad}

\section{Tareas individuales (entrega escrita)}

\begin{enumerate}[leftmargin=1.5em]
  \item \textbf{Tarea conceptual.}
  Explique, en máximo dos páginas, por qué la relación
  \begin{equation*}
    D_L=(1+z)^2D_A
  \end{equation*}
  es una pieza clave para conectar teoría y observación.

  \item \textbf{Tarea matemática.}
  Partiendo de
  \begin{equation*}
    \dot\rho + 3H\left(\rho+\frac{p}{c^2}\right)=0,
  \end{equation*}
  derive paso a paso $\rho\propto a^{-3(1+w)}$ para $p=w\rho c^2$.

  \item \textbf{Tarea numérica.}
  Usando un conjunto de datos sintético (mínimo 20 puntos) de $D_L(z)$,
  ajuste $H_0$ y $\OmM$ en un modelo plano materia+energía oscura.
  Reporte: función objetivo, parámetros óptimos, incertidumbres y gráfico de residuos.

  \item \textbf{Tarea crítica de literatura.}
  De los tres artículos base, elija uno y escriba una reseña técnica con esta estructura:
  problema, método, resultado principal, limitaciones, propuesta de mejora.

  \item \textbf{Tarea de extensión.}
  Proponga una mini-investigación para el semillero (2 a 4 semanas) inspirada en
  los artículos base. Debe incluir pregunta, método, datos necesarios y criterios de éxito.
\end{enumerate}

\section{Rúbrica sugerida}
\begin{itemize}[leftmargin=1.5em]
  \item precisión matemática y uso correcto de unidades: 30\%;
  \item claridad conceptual y argumentación física: 25\%;
  \item calidad del análisis de incertidumbre y validación: 20\%;
  \item lectura crítica de literatura científica: 15\%;
  \item presentación y comunicación científica: 10\%.
\end{itemize}

\chapter{Cierre}

Esta clase extensa ofrece una ruta completa para estudiar temas avanzados de cosmología
aplicada y análisis de datos astrofísicos. La combinación de teoría, matemática,
inferencia y lectura crítica permite que los estudiantes avancen desde fundamentos
hasta prácticas cercanas al trabajo real de investigación.

La recomendación final es iterar constantemente entre ecuaciones, datos y discusión
científica: calcular, contrastar, cuestionar y volver a calcular. Esa dinámica es la
base del pensamiento astronómico riguroso.

\backmatter
\bibliographystyle{apalike}
\bibliography{clase_cosmologia_articulos}

\end{document}
